%! Skills Focus: Selecting an Appropriate Inference Procedure for Categorical Data — Unit 8, Lesson 8.7 — Group Activity
\documentclass[10pt]{article}
\usepackage{apstats-article}
\usepackage{apstats-boxes}

\begin{document}

\pageheader{Unit 8, Lesson 8.7}{Group Activity}
\namepartnerperiod

\vspace{0.08in}

\textbf{Procedure Menu} (use for Tier R): one-sample $z$-interval for $p$ \;|\;
one-sample $z$-test for $p$ \;|\; two-sample $z$-test for $p_1-p_2$ \;|\;
chi-square goodness-of-fit test \;|\;
chi-square test for homogeneity \;|\;
chi-square test for independence

\vspace{0.06in}

\begin{scenariobox}[Scenario 1 --- Streaming Preferences]{sky}
A media company takes one random sample of 400 adults and records two things: each person's
age group (Under 35 or 35 and Over) and their primary streaming service (Service A, B, or C).
The company wants to test whether there is an association between age group and streaming
service preference.

\textbf{Tier R:} Name the procedure and state one key condition.
\begin{itemize}
  \item Procedure: \blank{8cm}
  \item Key condition: \blank{7cm}
\end{itemize}

\textbf{Tier A:}
\begin{enumerate}
  \item Procedure name: \blank{8cm}
  \item $H_0$: \writeline
  \item $H_a$: \writeline
  \item Key condition to check (state the requirement with numbers from the problem if possible):
        \writeline
\end{enumerate}
\end{scenariobox}

\begin{scenariobox}[Scenario 2 --- M\&M Color Distribution]{sky}
A statistics student purchases a large bag of M\&Ms. The manufacturer claims the color
distribution is 24\% blue, 20\% orange, 16\% green, 14\% yellow, 13\% red, and 13\% brown.
The student counts colors in a random sample of 200 M\&Ms from the bag and wants to test
whether her bag matches the claimed distribution.

\textbf{Tier R:} Name the procedure and state one key condition.
\begin{itemize}
  \item Procedure: \blank{8cm}
  \item Key condition: \blank{7cm}
\end{itemize}

\textbf{Tier A:}
\begin{enumerate}
  \item Procedure name: \blank{8cm}
  \item $H_0$: \writeline
  \item $H_a$: \writeline
  \item Minimum expected count for blue: $200 \times 0.24 =$ \blank{2cm} \quad Condition met? \blank{2cm}
\end{enumerate}
\end{scenariobox}

\begin{scenariobox}[Scenario 3 --- Vaccination Status and Opinion]{sky}
A public health researcher randomly samples 150 vaccinated adults and 150 unvaccinated adults.
She asks each person whether they support a new public health mandate (Support, Neutral, or
Oppose). She wants to test whether the distribution of opinion is the same for vaccinated and
unvaccinated adults.

\textbf{Tier R:} Name the procedure.
\begin{itemize}
  \item Procedure: \blank{8cm}
\end{itemize}

\textbf{Tier A:}
\begin{enumerate}
  \item Procedure name: \blank{8cm}
  \item $H_0$: \writeline
  \item $H_a$: \writeline
  \item Explain in one sentence why this is \emph{not} a test for independence: \writeline \writeline
\end{enumerate}
\end{scenariobox}

\begin{scenariobox}[Scenario 4 --- Voter Turnout]{sky}
A political researcher randomly selects 500 registered voters and records whether each voted
in the last election (yes or no). He wants to test whether the proportion who voted is
greater than 55\%.

\textbf{Tier A:}
\begin{enumerate}
  \item Procedure name: \blank{8cm}
  \item $H_0$: \blank{5cm} \quad $H_a$: \blank{5cm}
  \item Large Counts condition: $500 \times 0.55 =$ \blank{2cm} $\ge 10$? \blank{1.5cm} \quad
        $500 \times 0.45 =$ \blank{2cm} $\ge 10$? \blank{1.5cm}
\end{enumerate}
\end{scenariobox}

\begin{scenariobox}[Scenario 5 --- Flu Shot Rates]{sky}
A hospital system randomly samples 200 patients from its urban clinic and 200 patients from
its rural clinic. It records whether each patient received a flu shot this season (yes or no).
The hospital wants to test whether the proportion who received a flu shot differs between the
two clinics.

\textbf{Tier A:}
\begin{enumerate}
  \item Procedure name: \blank{8cm}
  \item $H_0$: \blank{5cm} \quad $H_a$: \blank{5cm}
  \item Why is this a two-sample $z$-test rather than a chi-square test for homogeneity?
        \writeline \writeline
\end{enumerate}
\end{scenariobox}

\begin{scenariobox}[Scenario 6 --- Tier E: Compare and Contrast]{sky}
\textbf{Tier E only.}

A researcher takes a random sample of 300 college students and records two variables: class
year (First-Year, Sophomore, Junior, Senior) and whether each student has a declared major
(yes or no).

\begin{enumerate}
  \item A student says this situation calls for a two-sample $z$-test for $p_1 - p_2$
        because the outcome (declared major) is binary. Identify the error in this reasoning
        and name the correct procedure.
        \writeline \writeline \writeline
  \item A second student says this calls for a chi-square goodness-of-fit test. Identify the
        error and explain why.
        \writeline \writeline \writeline
  \item State the correct $H_0$ and $H_a$ for the appropriate procedure.
        \begin{itemize}
          \item $H_0$: \writeline
          \item $H_a$: \writeline
        \end{itemize}
\end{enumerate}
\end{scenariobox}

\end{document}
